%!TEX root=../../main.tex

\limitingActivationSet*
\begin{proof}
    We say \( \Delta \preceq \Delta' \) if every entry of
    \( \Delta' \) is at least as large as \( \Delta \).
    If \( \Delta \prec \Delta' \), then it is straightforward to embed
    any network weights \( \theta \) conformal to \( \Delta \) into a network
    conformal \( \Delta' \) by setting the additional neurons to zero.
    Thus, \( \calDl_{\Delta} \subseteq \calDl_{\Delta'} \).
    Since
    \( \calDl_{\Delta} \subseteq \cbr{ \diag(z) : z \in \cbr{0, 1}^n} \),
    any monotone increasing sequence \( \Delta_{k} \) defines a bounded
    monotone increasing sequence \( \calDl_{\Delta_k} \) which must have
    a limit.
    If every entry of \( \Delta_k \) diverges as \( k \rightarrow \infty \),
    then this limit is the same no matter the choice of sequence.
    We call it \( \calDl \).
\end{proof}

\layerElimination*
\begin{proof}
    We start by defining a suitable \( T^{(l+1)} \) as follows:
    \begin{equation}\label{eq:t-plus-characterization}
        T^{(l+1)}_{j_1, \ldots, j_l}
        = \sum_{i_l \in \calI^{(l)}_{j_l}}
        T^{(l)}_{j_1, \ldots, j_{l-1}, i_l} \otimes W^{(l+1)}_{i_l}
        \text{ where } \calI^{(l)}_{j_{l}}
        = \cbr{ i_l : T^{(l)}_{i_l} \in \calK^{(l)}_{j_{l}}}
        ,
    \end{equation}
    where we recall that the index \( j_l \) ranges over the activation
    patterns achievable by the \( l^{\text{th}} \) layer,
    \[
        \calD^{(l)} = \cbr{
            D^{(l)}_{j_l}
            = \text{diag}\rbr{\mathds{1}(X^{(l)} \odot T^{(l)}_{i_l} \geq 0) }
            : T^{(l)} \in \calG^{(l)}, \, i_l \in [d_l]
        }.
    \]
    Remember that \( p_l = \abs{\calD^{(l)}} \) denotes the cardinality of this set.
    Using this notation, we can re-index the sum in
    \cref{eq:layer-to-tensorize} by activation patterns as follows:
    \begin{align*}
        Z^{(l+2)}
         & = \rbr{\sum_{i_l=1}^{d_l} \rbr{X^{(l)} \odot T^{(l)}_{i_l}}_+
        W^{(l+1)}_{i_l}}_+                                               \\
         & = \rbr{\sum_{j_l=1}^{p_l} \sum_{i_l \in \calI^{(l)}_{j_l}}
            \rbr{X^{(l)} \odot T^{(l)}_{i_l}}_+
            W^{(l+1)}_{i_l}}_+.
        \intertext{Linearizing the activation pattern now gives,}
        Z^{(l+2)}
         & = \rbr{\sum_{j_l=1}^{p_l}
            \sum_{i_l \in \calI^{(l)}_{j_l}}
            D^{(l)}_{j_l} \sbr{X^{(l)} \odot T^{(l)}_{i_l}}
        W^{(l+1)}_{m, \cdot}}_+                                          \\
         & = \rbr{\sum_{j_l=1}^{p_l}
        \sum_{j_1, \ldots, j_{l-1}} D^{(l)}_{j_l} X^{(l)}_{j_1, \ldots j_{l-1}}
        \rbr{
        \sum_{i_l \in \calI^{(l)}_{j_l}}
        T^{(l)}_{j_1, \ldots, j_{l-1}, i_l}
        \otimes W^{(l+1)}_{i_l}}}_+                                      \\
         & = \rbr{X^{(l+1)} \odot T^{(l+1)}}_+.
    \end{align*}
    It should be clear from the way that we constructed \( T^{(l+1)} \) that
    it is contained \( \calG^{(l+1)} \).
    This establishes one direction of theorem.
    The other direction can be shown by instead assuming that
    \( T^{(l+1)} \in \calG^{(l+1)} \) and then noting that the characterization
    in \cref{eq:t-plus-characterization} immediately follows for some
    \( T^{(l)} \in \calG^{(l)} \) and \( W^{(l+1)} \in \R^{d_l \times d_{l+1}} \).
    The preceding calculation is then reversed to obtain
    \cref{eq:layer-to-tensorize}.
\end{proof}

\tensorEquivalence*
\begin{proof}
    The argument is based on straightforward inductive application of
    \cref{lemma:tensor-elimination}.
    Let \( T^{(1)} = W^{(1)} \), which is in
    trivially in \( \calG^{(1)} = \R^{d_0 \times d_1} \).
    Thus, we have
    \[
        Z^{(2)} = \rbr{X^{(1)} \odot T^{(1)}}_+,
    \]
    which forms our base case.
    Now assume for \( l \geq 1 \) that
    \[
        Z^{(l+1)} = \rbr{X^{(l)} \odot T^{(l)}}_+,
    \]
    for some \( T^{(l)} \in \calG^{(l)} \) (or that \( l = k \), in which case
    there is no activation and we are done).
    By definition of the ReLU MLP,
    \[
        Z^{(l+2)} = \rbr{\rbr{X^{(l)} \odot T^{(l)}}_+ W^{(l+1)}}_+.
    \]
    \cref{lemma:tensor-elimination} now implies an equivalent tensorized
    version as follows:
    \[
        Z^{(l+2)} = \rbr{X^{(l+1)} \odot T^{(l+1)}}_+,
    \]
    for some \( T^{(l+1)} \in \calG^{(l+1)} \).
    Induction (with a minor modification to drop the ReLU in the final step)
    now implies there exists an equivalent
    tensor problem \( f_\theta(X) = X^{(k)} \odot T^{(k)} \) with
    \( T^{(k)} \in \calG^{(k)} \).
    Since \cref{lemma:tensor-elimination} establishes an if-and-only-if
    relationship, this argument is reversible to show that every
    tensor model corresponds to a non-convex ReLU network.
\end{proof}


\aside{Not ready for this result yet!}
This allows us to write an explicit equation for \( \calDl \) as follows.

\begin{restatable}{lemma}{activationCharacterization}\label{lemma:activation-characterization}
    The limiting set of activation patterns given in
    \cref{eq:limiting-activation-set} satisfies
    \begin{equation}\label{eq:explicit-activation-def}
        \calDl =
        \cbr{\Dl_{j_l} = \mathds{1}\rbr{\emph{\diag}
                \rbr{ \Xl \odot U^{(l)} \geq 0 }} :
            U^{(l)} \in \R^{d_1 \times p_1 \times \cdots \times p_{l}}}
    \end{equation}
\end{restatable}
\begin{proof}
    Define
    \[
        \tilde{\calD}^{(l)} =
        \cbr{\Dl_{j_l} = \mathds{1}\rbr{\diag
                \rbr{ \Xl \odot U^{(l)} \geq 0 }} :
            U_l \in \R^{d_0 \times p_1 \times \cdots \times p_{l-1}}}
    \]
    and let \( \Dl \in \calDl \) as defined in
    \cref{eq:limiting-activation-set}.
    By definition,
    \[
        \Dl_{j_l} =
        \diag\rbr{\mathds{1}\rbr{f_{\theta_{l-1}}(X) u \geq 0}},
    \]
    for some \( \theta_{l-1} \in \Theta^{(l-1)} \) and \( u \in \R^{d_l} \).
    Let \( U^{(l)} \in \R^{d_{l-1} \times d_{l}} \) such that
    \( U^{(l)}_{i_l} = u \).
    Then we can write
    \begin{align*}
        f_{\theta_{l-1}}(X) u
         & = \rbr{Z^{(l-1)} W^{(l-1)}}_+ U^{(l)}_{i_l}.
    \end{align*}
    \cref{lemma:layer-elimination} implies that there exists
    \( T^{(l)} \in \calG^{(l)} \) such that
    \[
        \rbr{Z^{(l-1)} W^{(l-1)}}_+ U^{(l)}
        = X^{(l)} \odot T^{(l)},
    \]
    and, in particular, that
    \[
        f_{\theta_{l-1}}(X) u
        = X^{(l)} \odot T^{(l)}_{i_l}.
    \]
    We deduce that \( \Dl_{j_l} \in \tilde{\calD}^{(l)} \) as required.

    \aside{TODO: We're not ready for this reverse direction.
        We need to consider wide networks and derive the limit of \( \calGl \)
        first.}
    Reversing this proof and using the only-if implication
    of \cref{lemma:layer-elimination} shows \( \tilde{\calD}^{(l)} \)
\end{proof}

